\documentclass[10pt]{article}
\usepackage[margin=0.78in]{geometry}
\usepackage{amsmath,amssymb,amsthm,mathtools}
\usepackage{graphicx,booktabs,array,xcolor,microtype,hyperref,placeins,multicol}
\hypersetup{colorlinks=true,linkcolor=blue!45!black,citecolor=blue!45!black,urlcolor=blue!55!black}
\setlength{\parindent}{0pt}
\setlength{\parskip}{3pt}
\setlength{\floatsep}{6pt plus 2pt minus 2pt}
\setlength{\textfloatsep}{8pt plus 2pt minus 2pt}
\setlength{\intextsep}{8pt plus 2pt minus 2pt}
\setlength{\abovecaptionskip}{4pt}
\setlength{\belowcaptionskip}{0pt}
\raggedbottom
\newcommand{\figfile}[1]{figures/#1.pdf}
\newcommand{\guidefig}[1]{figures/#1.pdf}
\newcommand{\verifyfig}[1]{figures/#1.png}
\setcounter{topnumber}{2}
\setcounter{bottomnumber}{1}
\setcounter{totalnumber}{3}
\renewcommand{\topfraction}{0.90}
\renewcommand{\bottomfraction}{0.80}
\renewcommand{\textfraction}{0.08}
\renewcommand{\floatpagefraction}{0.70}
\newcommand{\E}{\mathbb E}
\newcommand{\Prob}{\mathbb P}
\newcommand{\dd}{\,\mathrm d}
\newtheorem{theorem}{Theorem}
\newtheorem{lemma}{Lemma}
\newtheorem{corollary}{Corollary}

\title{Point-Trap Design in Diffusive Media: Minimax Capture, Graph Bounds, and Spectral Placement}
\author{Aneesh Arnav Chikkala\\\small October 2026}
\date{}

\begin{document}
\maketitle

\begin{abstract}
	In many transport settings, a diffusing particle is removed only when it encounters one of a limited number of reactive sites. The placement of these sites is therefore an optimization problem: with a fixed total amount of reactivity, the designer chooses their locations and strengths to control the longest capture time over all starting positions.
	
My study examines reflected diffusion in a heterogeneous one-dimensional medium with finitely many partially absorbing point traps. My main result reduces the continuum placement problem to a finite source-partition problem. A balancing principle equalizes resistance-weighted rises in each cell, and a monotone budget equation determines the strengths. The source share $p_i$, barrier $R_i$, and strength $\kappa_i$ satisfy $p_iR=\kappa_i(T-R_i)$, so source share is proportional to strength times remaining barrier capacity.
	
	I also derive a mass--resistance bound for metric graphs and a spectral result for long-time survival with prescribed strengths. In a saturated-sediment benchmark, the optimized placement lowers the modeled worst-start lifetime by $4.246\%$ relative to equal physical-depth placement. The interval theorem is formally checked in Lean under its stated assumptions, while the framework provides a testbed for related open problems, including narrow-capture and spectral-partition questions connected with the Hot Spots Conjecture. Together, these results provide a route from heterogeneous transport coefficients to optimized trap placement, with the sediment benchmark illustrating a possible application.
\end{abstract}

\section{Introduction and objective}

The central question is a point-trap design problem. A diffusing particle moves through a finite channel and is captured when it encounters one of finitely many localized partially absorbing sites. We may choose the trap locations
\[
0<x_1<\cdots <x_N<L
\]
and nonnegative strengths $\kappa_1,\ldots,\kappa_N$ subject to
\[
\sum_{j=1}^{N}\kappa_j=K.
\]
The design protects the least favorable starting point.

For a design $\mathcal D=(x_j,\kappa_j)_{j=1}^N$, let $\tau_{\mathcal D}$ be the first capture time and set
\[
u_{\mathcal D}(x)=\E_x[\tau_{\mathcal D}],
\qquad
M_N^*(K)=\inf_{\mathcal D:\,\sum\kappa_j=K}\;\max_{x\in[0,L]}u_{\mathcal D}(x).
\]
The inner maximum identifies the worst starting position; the outer infimum chooses the layout that minimizes that lifetime. The resulting objective is worst-start mean capture time.

Two companion objectives use the same designs but measure different observables. The spectral companion maximizes the principal decay rate $\lambda_1$ of the killed diffusion. A discounted first-capture question asks whether
\[
\sup_{\mathcal D}\;\inf_{x\in[0,L]}\E_x[e^{-s\tau_{\mathcal D}}]
\]
is attained by the same symmetric design. The interval theorem below concerns mean capture time; the spectral and discounted questions require their own arguments because a decay rate or discounted probability can rank layouts differently.

The interval reduction and graph equality conditions are proved. The discounted question remains a conjecture because its equation contains an additional zeroth-order term. The verification plots show sampled landscapes and estimator behavior; the universal statements come from the reductions and inequalities in Sections 3--4.

\begin{figure}[htbp]
\centering
\includegraphics[width=.68\linewidth]{\verifyfig{two_trap_landscape}}
\caption{Representative two-trap objective landscape from the verification results. The sampled values show the neighborhood of a numerical candidate.}
\label{fig:syn-landscape}
\end{figure}

The landscape shows the neighborhood of a two-trap candidate and its flat directions. The analytic bounds in Sections 3--4 are independent of this grid.

\begin{figure}[htbp]
\centering
\includegraphics[width=.68\linewidth]{\verifyfig{monte_carlo_convergence}}
\caption{Monte Carlo diagnostic with horizon $2$ and $6{,}000$ paths. The estimates rise from $0.25310$ to $0.25557$ as $\Delta t$ decreases from $0.002$ to $0.00025$; they remain below the finite-horizon target $0.2591433$ and full mean $0.2592593$. The gap records timestep, finite-band, and omitted-tail bias.}
\label{fig:syn-verification-convergence}
\end{figure}

Finite-band, timestep, and horizon-truncation biases remain visible in this estimator diagnostic.

\section{The channel in simple mathematical language}

The physical picture is a one-dimensional channel $x\in[0,L]$. A particle moves randomly along the channel, reflects at the two endpoints, and can be captured at the point sites $x_j$. The medium may be heterogeneous: its local transport coefficient and its local source density need not be constant. We write
\[
\mathcal L f(x)=\frac{1}{\rho(x)}\frac{\mathrm d}{\mathrm dx}
\left(a(x)\frac{\mathrm df}{\mathrm dx}\right),
\qquad a(x)>0,\quad \rho(x)>0.
\]
Here $a$ measures how readily motion transports across a small interval, while $\rho$ measures how much diffusive time is stored there. In physical terms, a region with large $\rho/a$ is slow to cross.

The traps are idealized as point interactions. The state continues through each $x_j$, but the derivative has a jump proportional to the trap strength:
\[
\bigl[a u'\bigr]_{x_j^-}^{x_j^+}=\kappa_j u(x_j).
\]
The reflecting endpoints satisfy $a(0)u'(0)=a(L)u'(L)=0$. For a fixed design, the mean first-capture time solves
\[
\mathcal L u=-1\quad\text{away from the traps}.
\]
The right side $-1$ is the clock: every unit of physical time contributes one unit to the expected lifetime. The jump condition is the local sink: a larger $\kappa_j$ removes probability more strongly when the particle reaches $x_j$.

The stochastic interpretation is the same statement in probabilistic language. The diffusion follows the reflecting dynamics generated by $\mathcal L$, and the traps kill it through local time:
\[
\Prob_x(\tau_{\mathcal D}>t)
=\E_x\!\left[\exp\left(-\sum_{j=1}^N\kappa_j L_t^{x_j}\right)\right],
\]
with the normalization of local time fixed by the jump condition. This is why a point trap is not a zero-width version of an arbitrary finite sink: its strength is a local-time coupling, and the limiting operator must retain that coupling \cite{Karatzas1991,Fukushima2011,Bressloff2022}.

The primary design problem is therefore
\[
\boxed{\;\text{choose }(x_j,\kappa_j)\text{ with }\sum_j\kappa_j=K
\text{ to minimize }\max_x u(x).\;}
\]
The principal eigenvalue problem,
\[
\mathcal L\phi=-\lambda\phi,
\qquad
\bigl[a\phi'\bigr]_{x_j^-}^{x_j^+}=\kappa_j\phi(x_j),
\]
is a second objective. Its largest admissible $\lambda=\lambda_1$ controls the slowest long-time decay, and the Rayleigh principle gives an efficient comparison tool. It is a companion analysis, not a redefinition of the mean-capture problem.

\section{Analytic minimax construction}

\subsection{Source coordinates and the physical meaning of a cell}

Define the normalized source coordinate
\[
q(x)=\frac{1}{R}\int_0^x\rho(s)\,\mathrm ds,
\qquad
R=\int_0^L\rho(s)\,\mathrm ds.
\]
The channel becomes $q\in[0,1]$. A source share $p$ means that a fraction $p$ of the total diffusive source mass lies in the corresponding physical interval. This coordinate is not a claim that the physical channel is uniform: the reconstruction
\[
\frac{\mathrm dx}{\mathrm dq}=\frac{R}{\rho(x(q))}
\]
returns the actual lengths after the optimization has been solved.

The coordinate change makes the conserved source measure transparent. In the transformed variable, the interior equation is controlled by a resistance density
\[
\omega(q)=\frac{R^2}{a(x(q))\rho(x(q))},
\]
while $\kappa_j$ remains the physical point strength. The proof can therefore use the normalized interval without discarding either the strength budget or the heterogeneous transport.
It is useful to distinguish the normalized budget $\widehat K:=K/R$ from the physical total strength $K$. The shares $p_i$ below are fractions of total source mass, whereas $\kappa_i$ remains a physical point strength.

\begin{figure}[htbp]
\centering
\includegraphics[width=.70\linewidth]{\guidefig{f03-source-coordinate}}
\caption{Source-coordinate normalization. The transformation changes the coordinate used by the proof while leaving the physical objective and strength budget explicit in the reconstruction.}
\label{fig:source-coordinate}
\end{figure}

\begin{figure}[htbp]
\centering
\includegraphics[width=.69\linewidth]{\guidefig{f05-balanced-rises}}
\caption{Balanced rises across a cell. The equal-rise condition is the local minimax step used before the global partition search.}
\label{fig:balanced-rises}
\end{figure}

\subsection{One-cell calculation and why equal-rise is forced}

The local step can be written without appealing to a picture. For a source-coordinate cell $[a,b]$, let $c$ be the trial trap location and let $\omega>0$ be the transformed resistance density. The two sides contribute
\[
A(c)=\int_a^c(c-s)\omega(s)\,\mathrm ds,
\qquad
B(c)=\int_c^b(s-c)\omega(s)\,\mathrm ds.
\]
The first quantity is the delay accumulated by material arriving from the left; the second is the corresponding delay from the right. Since
\[
A'(c)=\int_a^c\omega(s)\,\mathrm ds>0,
\qquad
B'(c)=-\int_c^b\omega(s)\,\mathrm ds<0,
\]
the difference $A(c)-B(c)$ is strictly increasing. There is therefore one and only one point $c_i$ at which the rises agree. Writing the common value as $R_i=A(c_i)=B(c_i)$ makes the local minimax statement explicit: moving the trap in either direction increases the larger side.

The cell source share $p_i$ and point strength $\kappa_i$ enter through the common barrier level $T$. After the two sides have been balanced, the cell equation reduces to
\[
p_iR=\kappa_i(T-R_i),
\qquad
\kappa_i(T)=\frac{Rp_i}{T-R_i},
\qquad T>\max_iR_i.
\]
This identity is the bridge between geometry and budget. The geometry determines $R_i$; the scalar $T$ determines how much strength each cell must receive to reach the same barrier height.

\begin{figure}[htbp]
\centering
\includegraphics[width=.70\linewidth]{\guidefig{f06-local-minimax}}
\caption{The local minimax mechanism. The two one-sided rises move in opposite directions as the trial trap position changes, so the maximum is minimized at their unique crossing.}
\label{fig:local-minimax}
\end{figure}
\FloatBarrier

\subsection{The scalar budget root}

After the cells are balanced, the total strength constraint determines the common barrier level. If $T$ denotes that level, the reconstructed strength in cell $i$ is $\kappa_i=\kappa_i(T)$, and the budget equation is
\[
\Phi(T):=\sum_{i=1}^N\frac{p_i}{T-R_i}=\widehat K,
\qquad\text{equivalently}\qquad
\sum_{i=1}^N\kappa_i(T)=K.
\]
Each summand in $\Phi$ is positive and strictly decreasing for $T>R_i$, so $\Phi$ is strictly decreasing on $T>\max_iR_i$. It diverges as $T$ approaches the largest active barrier from above and tends to zero as $T\to\infty$, giving a unique root for every positive physical budget $K$ (equivalently every positive $\widehat K$). The root can therefore be evaluated by certified bisection or a monotone-root routine rather than by an unconstrained nonlinear optimizer.

\FloatBarrier
\begin{figure}[htbp]
\centering
\includegraphics[width=.69\linewidth]{\guidefig{f09-budget-root}}
\caption{The monotone budget-root equation. The root is evaluated only after the cell barriers have been balanced and is bracketed above the largest barrier.}
\label{fig:budget-root}
\end{figure}

The construction now has a finite-dimensional outer problem: partition $[0,1]$ into source shares and solve one scalar root for each candidate partition. This is the point at which the original infinite-dimensional placement question becomes a controlled partition problem. The remaining search is still nontrivial, but its variables have a direct meaning and the local singular behavior has already been resolved.

\subsection{The theorem and proof logic}

Let $q_0=0<q_1<\cdots<q_{N-1}<q_N=1$ be a partition, with cell shares $p_i=q_i-q_{i-1}$. Let $t(q_1,\ldots,q_{N-1})$ be the monotone budget root obtained from the balanced-cell reconstruction. The central statement is:

\begin{theorem}[constructive minimax reduction]
For positive resistance density and positive total strength, every admissible point-trap design admits a balanced-cell representation with no larger worst-start mean. The optimal value is the minimum of the common barrier level over the compactified partition simplex, and the corresponding physical locations and strengths are recovered by the source-coordinate inverse and the monotone budget root.
\end{theorem}

The proof proceeds in four steps.

\begin{enumerate}
\item \emph{Decomposition.} Integrate the backward equation across each interval between adjacent traps. The flux is affine in the source mass, so the mean-time profile is a sum of cell contributions.
\item \emph{Local minimax.} If the two rises in a cell are unequal, moving the trap toward the larger rise lowers the larger one before increasing the smaller one to equality. Repeating this operation cannot increase the cell maximum; equality is therefore necessary at a minimizer.
\item \emph{Budget recovery.} Once the cell barriers are equalized, the strength allocation is the unique solution of the monotone equation $\Phi(t)=\widehat K$, equivalently $\sum_i\kappa_i(t)=K$. This removes a spurious degree of freedom from the placement search.
\item \emph{Attainment and converse.} The reduced objective extends continuously to the compactified partition simplex. A minimizing partition exists. Reversing the decomposition shows that every original design is bounded below by the reduced value, while the reconstructed design attains it.
\end{enumerate}

The global lower bound can be written as an explicit inequality chain. Let an arbitrary design $\mathcal D$ induce cell shares $p_i$, cell barriers $R_i$, site values $h_i=u_{\mathcal D}(x_i)$, and worst-start value $M=M(\mathcal D)$. The cell comparison and source balance give
\[
M-h_i\ge R_i,
\qquad
Rp_i=\kappa_i h_i,
\qquad
\frac{p_i}{M-R_i}\le\frac{p_i}{h_i}=\frac{\kappa_i}{R}.
\]
Summing over cells yields
\[
\Phi\bigl(M;q(\mathcal D)\bigr)
=\sum_i\frac{p_i}{M-R_i}
\le\frac{K}{R}=\widehat K.
\]
Since $\Phi(T(q);q)=\widehat K$ and $\Phi(\mathord\cdot;q)$ is strictly decreasing, this proves
\[
T\bigl(q(\mathcal D)\bigr)\le M(\mathcal D).
\]
Conversely, the equal-rise construction and the budget root reconstruct a physical design $\mathcal D(q)$ with $M(\mathcal D(q))=T(q)$. Minimizing over the compactified source-partition simplex gives
\[
\inf_{\mathcal D}M(\mathcal D)=\min_{q\in\Delta_{N-1}}T(q),
\]
where $\Delta_{N-1}$ is the compactified source-partition simplex.

\subsection{Sensitivity of the reduced objective}

The scalar root also gives an analytic sensitivity calculation that is useful when comparing partitions. Write
\[
w_i(T):=\frac{p_i}{(T-R_i)^2},
\qquad
S(T):=\sum_{i=1}^N w_i(T).
\]
Since $\Phi'(T)=-S(T)$, implicit differentiation of $\Phi(T)=K/R$ gives
\[
\frac{\partial T}{\partial K}=-\frac{1}{R S(T)},
\qquad
\frac{\partial T}{\partial R_j}=\frac{w_j(T)}{S(T)}.
\]
Thus additional total strength can only lower the common barrier, while increasing one cell's geometric barrier raises the global value in proportion to that cell's budget sensitivity. For a simultaneous perturbation of the cell barriers,
\[
\delta T
=\frac{\sum_i w_i(T)\,\delta R_i-\delta K/R}{S(T)}
+o\!\left(\lVert(\delta R,\delta K)\rVert\right).
\]
This formula separates the two sources of improvement: changing the partition reduces the weighted geometric term, while adding reactivity reduces the barrier directly.

The root has informative limiting regimes as well. Because $\sum_i p_i=1$,
\[
T=\frac{R}{K}+\sum_i p_iR_i+O(K/R),
\qquad K\downarrow0.
\]
At high strength, let $R_*:=\max_iR_i$ and $P_*:=\sum_{i:R_i=R_*}p_i$. Then
\[
T=R_*+\frac{RP_*}{K}+o(K^{-1}),
\qquad K\to\infty.
\]
The weak-reactivity regime is therefore controlled first by the source-weighted average barrier, whereas the strong-reactivity regime is controlled by the hardest active cell. These asymptotics provide analytic checks on the numerical partition search and explain why the maximum cell barrier cannot be ignored.

\begin{figure}[htbp]
\centering
\includegraphics[width=.70\linewidth]{\guidefig{f12-partition-landscape}}
\caption{Partition landscape used to distinguish an attained theorem from a local numerical candidate. The star marks a recorded numerical candidate; the analytic reduction specifies the admissible landscape being searched.}
\label{fig:partition-landscape}
\end{figure}

The figure shows the reduced objective and its boundary behavior. The proof is the pair of inequalities above.

\section{Graph and spectral companions}

\subsection{Weighted graph interpretation}

The same integrations work on a compact metric graph. An edge $e$ has length $\ell_e$, conductance $a_e$, and source density $\rho_e$; at a vertex $v$, the Kirchhoff balance includes the point reaction $J_v$. The natural graph operator is the weighted Neumann Laplacian
\[
Hf=-\frac{1}{\rho_e}(a_ef')'\quad\text{on each edge},
\qquad
\sum_{e\sim v}a_e\partial_\nu f(v)=J_v f(v).
\]
The graph variable $J_v$ is the same physical object as a point-trap strength: it is the amount of localized absorption assigned to that junction.

\begin{figure}[htbp]
\centering
\includegraphics[width=.68\linewidth]{\guidefig{f17-graph-model}}
\caption{Metric-graph model used for the conservation and deficit checks. Parallel edges, a loop, and a passive tip show why vertex incidence and edge weights must be tracked separately.}
\label{fig:graph-model}
\end{figure}

For a graph solution $Hu=1$, integration over all edges gives a conservation identity. Comparing the torsion profile with its edgewise linear profile produces three nonnegative deficits: reaction allocation, edge-profile curvature, and weighted mass--resistance imbalance. Equality requires the corresponding edge fluxes and weighted lengths to match.

Intervals, cycles, and reflecting stars have different vertex degrees but share the same equality condition: active gaps must be balanced after conductance and source weight are included.

\begin{figure}[htbp]
\centering
\includegraphics[width=.72\linewidth]{\guidefig{f24-equality-geometries}}
\caption{Equality geometries for the graph bound. The examples show how equal active gaps, rather than raw Euclidean distance alone, control equality in the weighted graph estimate.}
\label{fig:equality-geometries}
\end{figure}

\subsection{The graph deficit and the harmonic-mean obstruction}

Let $V=\sum_e\int_e\rho_e\,\mathrm dx$ be the total source mass, let $K=\sum_vJ_v$ be the total point strength, and let $C$ denote the effective conductance aggregate appearing in the edgewise estimate. The mean maximum satisfies
\[
M\;\ge\;\frac{V}{K}+\frac{V}{8C},
\]
and the difference from the right-hand side decomposes as
\[
\Delta_{\mathrm{graph}}
:=M-\frac{V}{K}-\frac{V}{8C}
=\Delta_{\mathrm{react}}+\Delta_{\mathrm{edge}}+\Delta_{\mathrm{mass}}\ge0.
\]
The first term vanishes when the reactive vertex values are balanced. The second vanishes when each edge profile has its sharp equality shape. The third records the arithmetic--harmonic obstruction: edge mass and resistance must match after vertex-degree and passive-tip factors are included.

The graph audit therefore reports both the total deficit and its component losses: a small total residual checks the identity, while equality requires each component to vanish.

\begin{figure}[htbp]
\centering
\includegraphics[width=.55\linewidth]{\guidefig{f22-deficit}}
\caption{Graph deficit components. The three nonnegative terms sum to the gap above the aggregate lower bound.}
\label{fig:graph-deficit}
\end{figure}

\subsection{Spectral placement and the sharp jump}

For the principal eigenvalue, the Rayleigh quotient is
\[
\lambda_1
=\inf_{f\ne0}
\frac{\displaystyle\int_0^L a|f'|^2\,\mathrm dx
      +\sum_j\kappa_j|f(x_j)|^2}
     {\displaystyle\int_0^L\rho|f|^2\,\mathrm dx}.
\]
The quotient makes the design tension explicit. A trap is useful where the ground-state mode has appreciable amplitude, but moving it changes the mode itself.

If $\phi>0$ is the normalized ground state, the comparison can be sharpened into a ground-state identity. With
\[
\mathcal E(v)=\int_0^L a|v'|^2\,\mathrm dx+\sum_j\kappa_j|v(x_j)|^2,
\qquad
\|v\|_\rho^2=\int_0^L\rho|v|^2\,\mathrm dx,
\]
one has, for every admissible $v$,
\[
\mathcal E(v)-\lambda_1\|v\|_\rho^2
=\int_0^L a\,\phi^2\left|\left(\frac{v}{\phi}\right)'\right|^2\,\mathrm dx\ge0.
\]
The right side is a sum of squares. Equality forces $v$ to be a scalar multiple of $\phi$. The finite collection of test functions in the validation figure checks the calculation numerically; the identity itself applies to the whole admissible energy space.

\begin{figure}[htbp]
\centering
\includegraphics[width=.55\linewidth]{\guidefig{f32-rayleigh-check}}
\caption{Rayleigh check for three independent test functions.}
\label{fig:rayleigh-check}
\end{figure}
\FloatBarrier

For one prescribed jump, write the incoming tangent as $t=\tan(\theta^-)$ and let the jump size be fixed. The phase update obeys the sharp one-jump inequality
\[
\arctan(t+2)-\arctan(t)\le 2\arctan(1),
\]
with equality at the centered phase. In physical terms, a fixed point strength removes the greatest possible phase advance when the incoming and outgoing modes are balanced around the site.

\begin{figure}[htbp]
\centering
\includegraphics[width=.70\linewidth]{\guidefig{f26-phase-jump}}
\caption{The sharp one-jump phase inequality behind prescribed-strength placement. The centered phase is the unique equality configuration for the fixed jump.}
\label{fig:phase-jump}
\end{figure}
\FloatBarrier

Iterating the jump inequality gives the placement theorem for prescribed strengths. Equal strengths give equal spacing in phase; a general strength vector is recovered from one monotone root and a map back to physical coordinates.

\FloatBarrier
\begin{figure}[htbp]
\centering
\includegraphics[width=.70\linewidth]{\guidefig{f28-spectral-root}}
\caption{Prescribed-strength phase root and reconstructed cell lengths. The root is solved in phase space; the lengths are then recovered in the physical coordinate.}
\label{fig:spectral-root}
\end{figure}

In the numerical comparison, strengths are held fixed and only locations change.

\begin{figure}[htbp]
\centering
\includegraphics[width=.70\linewidth]{\guidefig{f31-prescribed-comparison}}
\caption{Matched prescribed-strength comparison: phase sites versus midpoint sites. The phase construction changes only locations, making the spectral gain interpretable.}
\label{fig:prescribed-comparison}
\end{figure}

In the equal-strength case, the scalar root can be written with $h(\beta)$ and the parameter $\zeta$ satisfying
\[
\max\lambda_1=\frac{K}{L}\frac{\zeta}{\tan\zeta}<\frac{K}{L},
\qquad
\zeta\tan\zeta=\beta.
\]
The expansion
\[
\zeta\tan\zeta
=\zeta^2+\frac{\zeta^4}{3}+\frac{2\zeta^6}{15}+O(\zeta^8)
\]
provides a small-coupling check on the numerical root. It also shows why the strict inequality matters: diffusion between point sites leaves a finite spatial penalty even when total strength is fixed.

\section{Main solution, numerical values, and one application}

\subsection{What the theorem solves}

Table~\ref{tab:math-numbers} records three checks of the construction. The percentages compare the stated model outputs.

\begin{table}[htbp]
\centering
\scriptsize
\setlength{\tabcolsep}{2pt}
\caption{Straight numerical checks of the point-trap mathematics.}
\label{tab:math-numbers}
\begin{tabular}{>{\raggedright\arraybackslash}p{.25\linewidth}>{\raggedright\arraybackslash}p{.19\linewidth}>{\raggedright\arraybackslash}p{.25\linewidth}>{\raggedright\arraybackslash}p{.19\linewidth}}
\toprule
Check & Reference & Point-trap solution & Interpretation \\
\midrule
Heterogeneous one-trap mean & midpoint: $0.708333$ & $x=0.527525$, $M=0.688075$ & $2.860\%$ lower worst-start mean \\
Prescribed spectral strengths $(0.2,1.8)$ & midpoint $\lambda_1=1.651987$ & phase sites $(0.055965,0.555965)$, $\lambda_1=1.783518$ & $7.962\%$ higher principal rate \\
Homogeneous four-trap check & $D=L=1$, $K=4$ & sites $(.125,.375,.625,.875)$, $M=.255208$ & residual below $10^{-14}$ \\
\bottomrule
\end{tabular}
\end{table}

The heterogeneous one-trap row shows why geometric midpoint placement can miss the optimum. With strengths fixed, phase placement improves the slowest decay mode. The homogeneous four-trap row recovers the equally spaced symmetric limit.

\subsection{Real-world application of this optimization problem's solution}

Consider a one-dimensional saturated porous slab in which a dissolved hydrophobic pollutant diffuses through porewater and is sequestered by localized activated-carbon amendments. The carbon is a reactive sink in the reduced model, transferring pollutant from the mobile porewater phase into a carbon-bound reservoir. The point traps represent thin, highly reactive amendment bands, and the channel coordinate represents depth. This follows the transport-and-sequestration setting used in activated-carbon sediment and contaminant mass-transfer studies \cite{Ghosh2011,Werner2006}.

The reproducible benchmark sets $L=1$, $a(x)=1/(1+0.85x)$, $\rho(x)=1+0.55x$, and $K=2$, with two equal-width amendment bands of width $0.055$ and equal strengths. Equal physical-depth placement is the state-of-practice comparator. The finite-band reference calculation uses
\[
\rho c_t=(a c_x)_x-r(x)c,
\qquad
a c_x(0,t)=a c_x(L,t)=0,
\]
where $r(x)$ is supported on two narrow amendment bands. The theorem-guided layout uses the same total strength and band width, but places the bands at the reconstructed point-trap coordinates. The recorded sites are $0.250,0.750$ (equal depth), $0.384,0.711$ (equal source coordinate), $0.306,0.792$ (theorem-guided), and $0.296,0.789$ (finite-volume refinement); the last two are the layouts used in the headline comparison.

For the declared dimensionless slab and coefficient profile, the equal-depth comparator has worst-start lifetime $M_{\mathrm{base}}=0.722868$. The theorem-guided layout gives $M_{\mathrm{guided}}=0.692174$, a $4.246\%$ reduction. A direct finite-volume refinement gives $M_{\mathrm{refined}}=0.687172$, a $4.938\%$ reduction. The midpoint short-time capture diagnostic is $0.469327$ for the baseline, $0.470173$ for the guided layout, and $0.470197$ for the refined calculation. The modeled improvement is therefore a lower worst-case lifetime with almost no change in early capture.

The application literature motivates the transport setting, but the percentages still require calibration. Porous-media sorption, activated-carbon amendment, and contaminant sequestration add nonlinear sorption, advection, and finite-band effects beyond this benchmark. Mesocosm and field studies also show why bioavailability and water-column flux must be measured separately from a purely diffusive lifetime \cite{Tomaszewski2008,Beckingham2013}. The immediate use of the result is narrower: it provides a placement rule that can be tested before those mechanisms are added.

\textbf{Reading the application graph.} The gray, teal, and dashed orange curves are the equal-depth baseline, source-balanced design, and finite-volume refinement. Their worst-start lifetimes are $0.722868$, $0.692174$, and $0.687172$; the last two differ by $0.005002$.

\begin{figure}[htbp]
\centering
\includegraphics[width=.70\linewidth]{\figfile{fig03_state_of_practice_vs_point_trap}}
\caption{Application comparison: equal-depth placement, the theorem-guided point-trap design, and a direct finite-volume reference. The modeled reductions are $4.246\%$ and $4.938\%$.}
\label{fig:application-comparison}
\end{figure}
\FloatBarrier

Next tests should add advection, nonlinear sorption, and finite-width finite-rate amendment bands. A calibrated laboratory or field profile would allow the comparison to report uncertainty bands, mass removed, threshold time, and amendment cost alongside worst-start lifetime.

\section{Validation, limits, and the next conjecture}

The numerical record is specific. Refining the one-dimensional discretization from $101$ to $1601$ points changes the recorded value from $0.698317$ to $0.691805$, while the independent $801$-point reference is $0.692174$. Across the verification suite, the maximum linear-system residual over $5{,}000$ cases is $1.3378\times10^{-14}$; the maximum graph-deficit residual over $2{,}400$ cases is $3.6082\times10^{-16}$; and an independent finite-element comparison reports relative error $1.4275\times10^{-8}$. The scripts and fixed seeds are documented in \path{verification/README.md}; the JSON summaries are under \path{verification/results/} and \path{verification/generalized/results/}.

\subsection{What each validation layer actually tests}

Grid refinement tests stabilization, the residual sweep checks the solve after $A_hu_h=b_h$, the graph check tests the deficit identity, and the finite-element comparison uses another discretization for the same observable. The Monte Carlo diagnostic records a separate estimator: for $\Delta t=0.002,0.001,0.0005,0.00025$, its estimates rise from $0.25310$ to $0.25557$ while remaining below the finite-horizon target $0.2591433$ and full mean $0.2592593$. With horizon $2$ and $6{,}000$ paths, timestep, finite-band, and omitted-tail biases remain unresolved.

\pagebreak
The formal audit checks whether the recorded declarations type-check in the Lean kernel under the declared imports and foundational axioms. Numerical convergence and cross-solver agreement concern the computation; kernel checking concerns the encoded proof.

\subsection{Analytic meaning of the numerical checks}

The residual and convergence checks can be stated in operator form. If $A_h$ is the positive discrete operator, $b_h$ is the discretized source, and $\widetilde u_h$ is an approximate solution, define
\[
r_h=A_h\widetilde u_h-b_h.
\]
For the exact discrete solution $u_h=A_h^{-1}b_h$,
\[
u_h-\widetilde u_h=-A_h^{-1}r_h,
\qquad
\lVert u_h-\widetilde u_h\rVert_\infty
\leq \lVert A_h^{-1}\rVert_\infty\,\lVert r_h\rVert_\infty.
\]
The inverse-operator factor controls the residual-to-error estimate, while the second discretization probes consistency.

The mean-time observable has an equally useful integral representation. With
\[
S_{\mathcal D}(t,x):=\Prob_x(\tau_{\mathcal D}>t),
\]
Tonelli's theorem gives
\[
u_{\mathcal D}(x)=\E_x[\tau_{\mathcal D}]
=\int_0^\infty S_{\mathcal D}(t,x)\,\dd t.
\]
If the killed generator has eigenpairs $(\lambda_m,\psi_m)$ and the survival datum has coefficients $c_m(x)$, then formally
\[
S_{\mathcal D}(t,x)=\sum_{m\geq1}c_m(x)e^{-\lambda_mt},
\qquad
u_{\mathcal D}(x)=\sum_{m\geq1}\frac{c_m(x)}{\lambda_m}.
\]
Agreement in the principal eigenvalue alone is insufficient for the worst-start mean because the full weighted spectrum enters the integral.

The formal result is recorded in module \texttt{TimeTraps.Interval}: declarations \texttt{IntervalMedium}.\texttt{global\_minimum} and \texttt{IntervalDesign}.\texttt{optimal\_iff\_reconstruction}. The proof route is \path{formalization/review/INTERVAL.md}; the audit certificate is \path{formalization/audit/current-verification/result.json}. The archived environment is Lean 4.33.1 with Mathlib revision \texttt{0df444a3}; it records $2{,}282$ research modules, $10{,}543$ named theorems, $17{,}355$ audited declarations, and $47$ proved source groups, with only \texttt{propext}, \texttt{Classical.choice}, and \texttt{Quot.sound} as foundational axioms. Its scope is the encoded theorem and its stated assumptions \cite{Lean4,Mathlib}.

The next conjecture concerns a random deadline. Let $T_R$ be an independent exponential clock with rate $s$. Then
\[
\Prob_x(\tau<T_R)=\E_x[e^{-s\tau}],
\]
so the discounted objective has a direct interpretation as the probability of capture before an unpredictable disturbance. In the symmetric equal-strength family, the observed formula is
\[
\sup_{\mathcal D}\inf_x\E_x[e^{-s\tau_{\mathcal D}}]
\stackrel{?}{=}
\left[
\cosh\eta+
\frac{2N\sqrt{Ds}}{K}\sinh\eta
\right]^{-1},
\qquad
\eta=\frac{L}{2}\sqrt{\frac{s}{D}}.
\]
The phase inequality, symmetric calculation, and numerical sweep motivate this statement. In the recorded $N=6$ diagnostic, the largest discrepancy is $4.60\times10^{-4}$. The tested family does not include clustered, boundary, unequal-strength, or heterogeneous competitors, so the sweep is evidence for the conjecture rather than a proof.

If true in the full admissible class, the conjecture would extend the theory from expected lifetime to worst-start capture before a random deadline. The next proof target is a discounted equal-rise or phase inequality including the extra zeroth-order term. Nonlinear sorption, advection, open boundaries, and finite-width reactions are outside the current model.

\section*{Conclusion}

On the interval, the source-coordinate reduction gives the minimax layout and its physical reconstruction. The graph and spectral calculations provide separate extensions, and the numerical/formal records make the reported implementation auditable. The porous-media example remains a dimensionless screening benchmark; calibration and the discounted-capture conjecture are the next steps.

\begin{multicols}{2}
\begin{thebibliography}{99}
\scriptsize
\bibitem{Karatzas1991} I. Karatzas and S. E. Shreve, \emph{Brownian Motion and Stochastic Calculus}, 2nd ed., Springer, New York, 1991.
\bibitem{Fukushima2011} M. Fukushima, Y. Oshima, and M. Takeda, \emph{Dirichlet Forms and Symmetric Markov Processes}, 2nd ed., De Gruyter, Berlin, 2011.
\bibitem{Fitzsimmons1999} P. J. Fitzsimmons and J. Pitman, "Kac's moment formula and the Feynman--Kac formula for additive functionals of a Markov process," \emph{Stochastic Processes and their Applications} 79 (1999), 117--134.
\bibitem{Bressloff2022} P. C. Bressloff, "Diffusion-mediated absorption by partially-reactive targets: Brownian functionals and generalized propagators," \emph{Journal of Physics A: Mathematical and Theoretical} 55 (2022), 205001.
\bibitem{Albeverio1988} S. Albeverio, F. Gesztesy, R. H{\o}egh-Krohn, and H. Holden, \emph{Solvable Models in Quantum Mechanics}, Springer, Berlin, 1988.
\bibitem{Buttazzo2006} G. Buttazzo, N. Varchon, and H. Zoubairi, "Optimal measures for elliptic problems," \emph{Annali di Matematica Pura ed Applicata} 185 (2006), 207--221.
\bibitem{Tilli2020} P. Tilli and D. Zucco, "Optimal partitioning of an interval and applications to Sturm--Liouville eigenvalues," \emph{Journal of Differential Equations} 268 (2020), 4900--4919.
\bibitem{KurasovSerio2019} P. Kurasov and A. Serio, "Optimal potentials for quantum graphs," \emph{Annales Henri Poincar{\'e}} 20 (2019), 1517--1542.
\bibitem{Mugnolo2023} D. Mugnolo and M. Pl{\"u}mer, "On torsional rigidity and ground-state energy of compact quantum graphs," \emph{Calculus of Variations and Partial Differential Equations} 62 (2023), 27.
\bibitem{Lean4} L. de Moura and S. Ullrich, "The Lean 4 theorem prover and programming language," in \emph{Automated Deduction -- CADE 28}, Lecture Notes in Computer Science 12699, Springer, 2021, 625--635.
\bibitem{Mathlib} The mathlib Community, "The Lean mathematical library," in \emph{Proceedings of the 9th ACM SIGPLAN International Conference on Certified Programs and Proofs}, ACM, 2020, 367--381.
\bibitem{choi2014} Y. Choi et al., "Engineered sorbents for contaminant sequestration in porous media," \emph{Environmental Science and Technology} (2014).
\bibitem{qiu2023} Z. Qiu et al., "Activated carbon amendment for contaminant immobilization in sediment," \emph{Science of the Total Environment} (2023).
\bibitem{wikstrom2024} M. Wikstr{\"o}m et al., "Carbon amendments and contaminant transport in porous materials," \emph{Journal of Hazardous Materials} (2024).
\bibitem{pcb2006} United States Environmental Protection Agency, \emph{A Guide to Developing and Documenting Cost Estimates During the Feasibility Study}, EPA 540-R-00-002, 2006.
\bibitem{Ghosh2011} U. Ghosh, R. G. Luthy, G. Cornelissen, D. Werner, and C. A. Menzie, "In-situ sorbent amendments: A new direction in contaminated sediment management," \emph{Environmental Science \& Technology} 45(4) (2011), 1163--1168. DOI: 10.1021/es102694h.
\bibitem{Werner2006} D. Werner, U. Ghosh, and R. G. Luthy, "Modeling polychlorinated biphenyl mass transfer after amendment of contaminated sediment with activated carbon," \emph{Environmental Science \& Technology} 40(13) (2006), 4211--4218. DOI: 10.1021/es052215k.
\bibitem{Tomaszewski2008} J. E. Tomaszewski, P. B. McLeod, and R. G. Luthy, "Measuring and modeling reduction of DDT availability to the water column and mussels following activated carbon amendment of contaminated sediment," \emph{Water Research} 42(16) (2008), 4348--4356. DOI: 10.1016/j.watres.2008.07.016.
\bibitem{Beckingham2013} B. Beckingham and U. Ghosh, "Polyoxymethylene passive samplers to monitor changes in bioavailability and flux of PCBs after activated carbon amendment to sediment in the field," \emph{Chemosphere} 91(10) (2013), 1401--1407. DOI: 10.1016/j.chemosphere.2012.12.074.
\bibitem{PointTrapArchive} A. A. Chikkala, \emph{Point-Trap Design in Diffusive Media}, research archive containing the journal paper, Lean code, and handwritten mathematics, GitHub (2026). \url{https://github.com/aneesharnavch/Point-Trap-Design-in-Diffusive-Media}
\end{thebibliography}
\end{multicols}

\end{document}
